\section{Introduction}\label{sec:intro}

Loop fusion avoids repeated traversals and materialised intermediates~\citep{halide,tvm}.
However, when a later loop reads the final result of an earlier reduction, fusing the loops by interleaving their iterations would violate this dependency.
Such dependencies arise in tensor operators for machine learning and scientific computing.

% Shared unchanged by both prose editions.
\begin{wrapfigure}{r}{0.33\textwidth}
\vspace{-6pt}
\centering
\captionsetup{skip=3pt}
\begin{tikzpicture}[x=1cm,y=1cm,
 instance/.style={draw=origblue,rounded corners=2pt,fill=boxgray,inner sep=2pt,font=\scriptsize},
 carry/.style={->,draw=rulegray},
 dependency/.style={->,draw=staleorange,semithick}]
\node[instance] (a) at (0,0) {$I_1(0)$};
\node[instance] (b) at (1.15,0) {$I_1(1)$};
\node[font=\scriptsize,inner sep=1pt] (gap1) at (2.10,0) {$\cdots$};
\node[instance] (c) at (3.15,0) {$I_1(N-1)$};
\node[instance] (d) at (0,-0.95) {$I_2(0)$};
\node[instance] (e) at (1.15,-0.95) {$I_2(1)$};
\node[font=\scriptsize,inner sep=1pt] (gap2) at (2.10,-0.95) {$\cdots$};
\node[instance] (f) at (3.15,-0.95) {$I_2(N-1)$};
\node[font=\scriptsize,anchor=south west,inner sep=0pt]
  at ([yshift=2pt]a.north west) {Loop 1};
\node[font=\scriptsize,anchor=south west,inner sep=0pt]
  at ([yshift=2pt]d.north west) {Loop 2};
\foreach \s/\t in {a/b,b/gap1,gap1/c,d/e,e/gap2,gap2/f} \draw[carry] (\s) -- (\t);
\foreach \t in {d,e,f} \draw[dependency] (c.south) -- (\t.north);
\end{tikzpicture}
\caption{Schematic of dependent reductions.}
\label{fig:dependent-loops}
\vspace{-6pt}
\end{wrapfigure}

In a \emph{dependent reduction}, the reduction body uses a value determined by an earlier reduction's final result over the same axis.
Such a value is its \emph{context}: it is invariant along the reduction axis, but becomes available only after the earlier computation completes.
In \cref{fig:dependent-loops}, $I_i(j)$ denotes iteration $j$ of loop $i$, and the orange arrows show the dependency between the loops.

Fusing these reductions into a single pass therefore requires accumulating information for the later reduction before its final context is available.
An aggregate quantity retained for a part of the input is a \emph{state component}, or \emph{state}; the retained quantities form its \emph{state tuple}.

A \emph{join} combines the tuples of two disjoint parts into the tuple for their union.
For results computed with local contexts, this operation must account for the change from each part's context to that of the union.
The two connected questions are therefore which states retain enough information and how to derive a join that updates those states correctly.

Developing such algorithms requires specialised expertise.
Online normalisation, moment accumulation and tiled attention require identifying retained quantities, deriving updates, and implementing and tuning kernels~\citep{milakov,chan,pebay,flashattention}.
Recent automatic fusion methods handle algebraic corrections such as shifting or rescaling partial results~\citep{neptune,redfuser}.
Their correction rules do not cover dependencies that require aggregate information beyond the existing reduction results.

\textbf{Our Work.}
We present \sys,\footnote{Greek \emph{hapax}, ``once''.} which fuses dependent tensor reductions by augmenting source reduction results and deriving their joins.
We draw on methods developed for divide-and-conquer parallelisation: Farzan and Nicolet synthesise auxiliary accumulators and joins~\citep{farzan17,farzan21}, while AutoLifter synthesises auxiliary values and combinators for divide-and-conquer computations~\citep{autolifter}.
These methods introduce extra information so that results from independently processed parts can be combined.
\sys adapts this idea to fusion: \emph{auxiliary states} supplement source results with information needed for their context corrections.

\par\medskip
\noindent\begin{minipage}{\linewidth}
\begin{minipage}[t]{0.50\linewidth}
\vspace{0pt}
Each part computes its state tuple using its own contexts; a join corrects and combines these tuples (\cref{fig:dc-fusion}).
Although computing a larger part's tuple may take multiple traversals, a singleton tuple depends on just one input element.
In this formulation, a {\bfseries fused loop} can be viewed as sequentially joining singleton state tuples in one traversal.
\end{minipage}\hfill
% Shared English vector figure; both prose editions include this source.
% The introducing prose occupies the other half of the enclosing minipage.
\begin{paperdiagram}
\captionsetup{type=figure,skip=4pt}
\centering
\begin{tikzpicture}[x=1cm,y=1cm,>={Stealth[length=1.2mm,width=0.9mm]},
  label/.style={font=\footnotesize,align=center},
  instance/.style={draw=origblue,rounded corners=2pt,fill=boxgray,
    inner sep=2pt,minimum height=0.32cm,font=\footnotesize},
  continuation/.style={font=\footnotesize,inner sep=1pt},
  part/.style={draw=origblue,rounded corners=4pt},
  joinop/.style={draw=fixgreen,rounded corners=2pt,fill=fixgreen!8,
    inner xsep=6pt,inner ysep=3pt,font=\footnotesize\bfseries},
  carry/.style={->,draw=rulegray},
  dependency/.style={->,draw=staleorange,semithick},
  combine/.style={->,draw=fixgreen,semithick}]

\node[label] at (1.50,0) {Partition $T$};
\node[label] at (4.90,0) {Partition $S$};
\draw[part] (0,-0.22) rectangle (3.00,-2.10);
\draw[part] (3.40,-0.22) rectangle (6.40,-2.10);

% Symbolic endpoints stand for the first and last iterations of each part.
\node[instance] (t1a) at (0.55,-0.55) {$I_1(a)$};
\node[continuation] (t1dots) at (1.50,-0.55) {$\cdots$};
\node[instance] (t1b) at (2.45,-0.55) {$I_1(b)$};
\node[instance] (t2a) at (0.55,-1.50) {$I_2(a)$};
\node[continuation] (t2dots) at (1.50,-1.50) {$\cdots$};
\node[instance] (t2b) at (2.45,-1.50) {$I_2(b)$};
\node[instance] (s1c) at (3.95,-0.55) {$I_1(c)$};
\node[continuation] (s1dots) at (4.90,-0.55) {$\cdots$};
\node[instance] (s1d) at (5.85,-0.55) {$I_1(d)$};
\node[instance] (s2c) at (3.95,-1.50) {$I_2(c)$};
\node[continuation] (s2dots) at (4.90,-1.50) {$\cdots$};
\node[instance] (s2d) at (5.85,-1.50) {$I_2(d)$};
\foreach \a/\dots/\b in {t1a/t1dots/t1b,t2a/t2dots/t2b,s1c/s1dots/s1d,s2c/s2dots/s2d} {
  \draw[carry] (\a.east) -- (\dots.west);
  \draw[carry] (\dots.east) -- (\b.west);
}
% Each part's completed first reduction supplies its second reduction.
\foreach \a/\b in {t1b/t2a,t1b/t2b,s1d/s2c,s1d/s2d}
  \draw[dependency] (\a.south) -- (\b.north);
\node[label,text=auxpurple] at (1.50,-1.90) {Supplemental states};
\node[label,text=auxpurple] at (4.90,-1.90) {Supplemental states};

\node[joinop] (combined) at (3.20,-2.70) {Join};
\draw[combine] (1.50,-2.10) -- (combined.north west);
\draw[combine] (4.90,-2.10) -- (combined.north east);
\end{tikzpicture}
\caption{Partitioned computation.}
\label{fig:dc-fusion}
\end{paperdiagram}

\end{minipage}\par\medskip

To determine what information can support such a join, \sys separates reduction bodies from their contexts algebraically.
The resulting \emph{reference states} are reductions whose bodies do not read those contexts, hence are \emph{context-independent}.
They combine independently and determine the required results, establishing a single-pass construction when the analysis succeeds.
Their number gives a sufficient state-count bound independent of the reduction length.

The reference construction provides a valid single-pass representation, but executing its expanded arithmetic may undo source centring or scaling and introduce cancellation or overflow~\citep{chan,blanchard,boldo_fp}.
These numerical concerns motivate starting the \emph{maintained state tuple}---the quantities accumulated by the fused program---from source reduction results computed in local contexts.
Retaining these results keeps the source's centred or scaled quantities available for context corrections.
When the source results do not contain enough information for these corrections, auxiliary states supply the missing information.

\sys synthesises auxiliary states that, together with the source reduction results, form a maintained tuple convertible to and from the reference tuple.
For each candidate tuple, the synthesis procedure constructs and verifies the conversion expressions.
These conversions yield a correct join by recovering the reference tuples of the input parts, applying their known join, and reconstructing the maintained tuple of the union.
Thus, synthesis searches for auxiliary states satisfying this conversion requirement, without separately searching for join expressions.

To constrain the auxiliary-state search space, we construct candidate expressions from how the source reduction results respond to context changes.
Symbolically shifting or scaling a local context exposes aggregate coefficients needed to correct its result.
These coefficients provide candidate expressions for syntax-guided synthesis, restricting the search to quantities with a role in the update.
Together, fixed source components and these candidates limit what the synthesiser must invent.
The same corrections provide execution formulas that adjust local results while reusing source centring and scaling where possible.
Verifying these formulas against the join obtained through recovery connects the numerical formulation to the algebraic correctness argument.

Over the reals, the resulting join computes the maintained state tuple of the combined input.
Because the join also applies to arbitrary disjoint non-empty parts, the derived reducer admits independent computations on parts followed by a reduction tree.
Fewer traversals create opportunities for reducing data movement; realised performance also depends on state size, join cost and scheduling.

We implement \sys as a rewrite of TVM tensor expressions (TE)~\citep{tvm}.
A recorded implementation snapshot derives single-pass reducers for \nShowFused of \nShowcase selected operators from machine learning and scientific computing.
We report coverage, discovered states and derivation cost, and describe the implemented MPFR validation used to compare the emitted programs' accuracy with their sources.
GPU scheduling and performance evaluation remain subsequent work.

\noindent This paper makes the following contributions:
\begin{itemize}
\item \textbf{Fusion of dependent-reduction patterns.}
Context separation identifies a class of dependent chains with a single-pass construction, including consumer bodies outside the stated algebraic conditions of recent automatic fusion methods, and bounds the states sufficient for this construction.
\item \textbf{Reference-guided state synthesis and join derivation.}
Complete recovery of a reference representation reduces synthesis to state augmentation and makes a separate search for join expressions unnecessary.
\item \textbf{Source-guided search and numerical formulation.}
Source results and their context corrections restrict auxiliary-state search while providing execution formulas that reuse source centring and scaling where possible.
\item \textbf{Implementation and evaluation.}
A TVM implementation and recorded results on coverage, discovered states and derivation cost for \nShowcase operators.
\end{itemize}
